\section{Introduction}
\label{sec:intro}
Speculative decoding promises lower serial decoding latency by using a cheaper draft model to propose tokens that a target model validates together \citep{leviathan2023}. High acceptance, fast drafting, and cheap batched target execution can each make an implementation look promising. We ask what prevents those local indicators from translating into gains in integrated post-prefill speculative decode on the tested consumer configurations, and whether the measured constraints can guide runtime decisions.

We investigate MacBook Air M4, Core Ultra 9 275HX with RTX 5070 Ti Laptop, Core i5-13400 with Arc B570, and Core Ultra X7 358H with Arc B390 and CPU/NPU drafting. Outcomes differ sharply even within a platform. The representative M4 workload is token-exact but reaches only \MRepSpeed$\times$ speedup; a separately preregistered HumanEval workload on the same M4 is exact and reaches \MHEspeed$\times$. Differing session conditions prevent attributing this contrast to workload alone. Profitability is not a simple hardware capability in these studies.

Two diagnoses explain important failures elsewhere. On RTX 5070 Ti Laptop, correctness qualification removes cheap batching; the tested exact serialized controls cost more. B570 serialization improves agreement without qualifying the full workload. On 358H, the isolated NPU draft advantage does not survive the tested stateful integration, while complete target-side work dominates CPU drafting. These diagnoses neither exhaust possible causes nor explain the M4 contrast. Hardware, workload, and runtime co-vary; all gains concern post-prefill decode, excluding loading and initial prompt prefill.

Our contribution combines retrospective consumer-stack characterization with separately frozen runtime studies. The characterization retains positive, negative, and correctness-blocked outcomes, using path controls and direct timing to distinguish component promise from realized cost. Runtime qualification then blocks the tested 358H intervention before formal performance evaluation. On an exact M4 path, a guard using already-realized cost achieves profitable held-out execution against target-only, although fixed speculation has a higher aggregate point estimate. The studies connect correctness-qualified paths to measured control decisions. Performance characterization, heterogeneous placement, cost/commit accounting, and deterministic speculation have prior art \citep{dovetail,duodecoding,illusion,specllm,llm42}; our contribution is this bounded empirical connection, not a new speculative algorithm or universal profitability theory.
